\subsection{Related work}
\label{sec:related}

Multi-stage inspection has a reservation-index theory. \citet[Appendix~G]{kleinberg2016descending}
give generalized strike prices and \citet[Section~4, Definition~4]{bowers2024nested} a
constructive nested-index and capped-value recursion, both for terminal-only collection;
the latter also explicitly allow later costs to depend on preceding signals
\citep[Section~4.1]{bowers2024nested}.
\citet{gupta2019markovian} treat independent Markov inspection processes with rewards at
destination states, nonnegative prices and unbounded depth, of which the fixed-depth
positive-cost case is recovered here; \citet{boodaghians2023order} handle independent rewards
under line or tree precedence. Closest for intermediate collection is
\citet[Theorem~4.4]{yang2025markovian}: each reward on a Markov line is collectable, probing
costs are fixed, and the optimal policy selects the line with largest generalized reservation
value. With positive costs that is the corresponding special case here; their finite-support
dynamic program and our history-state covered-call construction are different
representations. The superprocess literature supplies sufficient conditions for an index
theorem \citep{whittle1980gittins,glazebrook1982whittle,scully2025gittins};
Appendix~\ref{app:remarks} identifies that condition for this model. More general Markov
search, in which a process offers a choice of actions, need not admit an index policy
\citep{glazebrook1982whittle,scully2025gittins}; one-shot search with an optional partial
inspection has none and is NP-hard \citep{aouad2025sequential}, as is combinatorial Markov
search \citep{bowers2026cms}.

Correlated search is also established: learning and approximation from joint-distribution
samples \citep{chawla2020correlations}, the reduction of fully adaptive correlated search to
Uniform Decision Tree \citep{chawla2023adaptivecorrelations}, correlation-aware reservation
rules \citep{gergatsouli2023weitzman}, and a tight factor four against the best partially
adaptive policy for a minimization objective \citep{bansal2025correlated}. Those approximation
factors do not transfer to our reward-minus-cost objective, and the correlated-BO stopping
guarantee of \citet{xie2025coststop}, no worse than stopping after one evaluation, is a
different kind of statement from Bellman optimality. In Bayesian optimization, PBGI is the
single-stage predecessor \citep{xie2024pbgi}; freeze-thaw and fidelity-aware surrogates
\citep{swersky2014,klein2017fabolas} motivate staged observation models; trace-aware knowledge
gradient \citep{wu2019takg}, successive halving and Hyperband
\citep{jamieson2016sh,li2018hyperband} allocate resources by other means. For this utility in
fixed-pool freeze-thaw optimization, \citet{vu2026coststop} analyze the one-step rule, expected
improvement stopping, and a certified two-step correction; the present paper gives the exact
index for the independent-arm staged problem.

\paragraph{Organization.}
Section~\ref{sec:setting} defines the model. Section~\ref{sec:main} states the
covered-call recursion, the index theorem and the loss decomposition, with
proofs; Sections~\ref{sec:comparative} and~\ref{sec:robust} treat comparative
statics and information sharing. Section~\ref{sec:validation} summarizes the
exact numerical verification, whose retained models, checks and results are in
Appendix~\ref{app:experiments}, and Section~\ref{sec:limitations} concludes.
Algorithm~\ref{alg:staged-index} gives the policy implementation in the main text.
